

Since our experiments use a finite $M$,
$\gcovbayeshat$ and $\gcovijhat$ are MCMC estimators of $\gcovbayes$ and
$\gcovij$, respectively.  Similarly, sampling with replacement from the rows of
$\xvec$ in \algrref{boot} is equivalent to drawing integer weights $\wvec \sim
\mathrm{Multinomial}(N; N^{-1}, \ldots, N^{-1})$ and computing the weighted
posterior expectation $\expect{\postw}{\g(\theta)}$ of \eqref{weighted_post}.
We therefore define the exact bootstrap covariance as the expectation under
multinomial resampling,
%
\begin{align*}
%
\gcovboot :={}& N \cov{\wvec \sim \mathrm{Multinomial}(N; N^{-1}, \ldots, N^{-1})}
    {\expect{\postw}{\g(\theta)}},
%
\end{align*}
%
of which $\gcovboothat$, computed with $B$ bootstrap samples and $M$ MCMC draws
for each, is a ``doubly Monte-Carlo'' estimator.  In general, for $B < \infty$,
we expect bootstrap sampling variability to remain even as $M \rightarrow
\infty$. Similarly, for $M < \infty$, we expect MCMC sampling variability to
remain even as $B \rightarrow \infty$.  We will shortly discuss how to estimate
these Monte Carlo errors.

Even when the conditions of \thmref{ij_consistent} are met and the IJ covariance
is consistent, we do not expect $\gcovij = \gcovboot$, even though $\gcovboot$ involves no Monte
Carlo error, since we have a finite number $N$ of data points, and the
bootstrap and IJ covariance are formally distinct (though each consistent)
estimators of $\gcovtrue$.  The difference $\gcovij - \gcovboot$ should decrease
as $N$ grows, but, for any finite $N$, the difference can be expected to exceed
Monte Carlo error for sufficiently large $M$ and $B$. Similarly, even in the
correctly specified case, we do not expect $\gcovbayes = \gcovij$, as, again,
these are distinct (though consistent) estimators of $\ggrad{1}(\thetatrue)
\info^{-1} \ggrad{1}(\thetatrue)^\trans$.

One might attempt to additionally account for frequentist sampling variability
of the original data from the target population in $\gcovij$ or $\gcovboot$ as
estimators of $\gcovtrue$, or of $\gcovbayes$ as an estimator of
$\ggrad{1}(\thetatrue) \info^{-1} \ggrad{1}(\thetatrue)^\trans$.
However, we expect the quantities $\gcovij$, $\gcovboot$, and $\gcovbayes$ to be
highly correlated under random sampling of the data from the population
distribution, and it is necessary to take this correlation into account when
computing the variability of, say, the IJ ``error'' $\gcovijhat - \gcovboothat$.
The authors do not see a simple, computationally tractable way to account for
this correlation, so we neglect this error in our experiments, and account only
for the Monte Carlo error of the preceding paragraph.  Since we consider the IJ
covariance  to succeed when the difference $\gcovijhat - \gcovboothat$ is small
relative to random variability, neglecting the frequentist sampling error is
conservative.

To estimate the Monte Carlo standard errors, we model the individual entries of
the matrices $\gcovbayeshat$, $\gcovijhat$, and $\gcovboothat$ as normal
distributions with Monte Carlo standard deviations given by the matrices
$\sebayes$, $\seij$, and $\seboot$, respectively.\footnote{Since both
$\gcovbayeshat$ and $\gcovijhat$ are based on the same MCMC samples, their Monte
Carlo errors are correlated, typically positively.  However, to avoid cumbersome
exposition of the experimental results, we will report standard errors for
$\gcovbayeshat$ and $\gcovijhat$ as if they were independent, since accounting
for this correlation will tend to reduce the standard errors of the difference
$\gcovbayeshat - \gcovijhat$, and the difference $\gcovbayeshat - \gcovijhat$
tends to be large in our experiments, even relative to the conservative estimate
of independent standard errors.
% Estimated standard errors of the difference
% $\gcovbayeshat - \gcovijhat$ taking this correlation into account are available
% in the online code for our experiments CITE GIT REPO.
}
%
To estimate $\sebayes$ and $\seij$, we block bootstrap the MCMC chains using
block sizes that are large relative to estimates of the MCMC mixing times.
Note that block bootstrapping here requires only re-arranging the original MCMC
chains, which is fast, and does not require drawing any new MCMC samples.  To
estimate $\seboot$, we use the fact that each bootstrap sample $\mathscr{G}^b$
from \algrref{boot} contains independent MCMC errors, and so the marginal
distribution of the $\mathscr{G}^b$ can be used to estimate the variability both
of the bootstrap sampling and the MCMC.  We do so using the delta method, as
follows.  For any two scalar random variables, $t_1$ and $t_2$ drawn according
to a joint distribution $p(t_1, t_2)$, we can define $h(a, b, c) :={} a - bc$
and write
%
\begin{align*}
%
\cov{p(t_1, t_2)}{t_1, t_2} &={}
    \expect{p(t_1, t_2)}{t_1 t_2} -
    \expect{p(t_1, t_2)}{t_1}\expect{p(t_1, t_2)}{t_2}\\
&={} h\left(\expect{p(t_1, t_2)}{t_1 t_2},
     \expect{p(t_1, t_2)}{t_1},
     \expect{p(t_1, t_2)}{t_2} \right).
%
\end{align*}
%
By applying the delta method to the function $h$, we can estimate the sampling
variability of $\cov{p(t_1, t_2)}{t_1, t_2}$ from the covariance of the vector
$(t_1 t_2, t_1, t_2)$, which, in turn, can be estimated from draws of $t_1$ and
$t_2$.  We apply the delta method in this way to each entry of the matrix
$\gcovboothat$.
